\section{Out-of-sample predictions}\label{sec:predictions}

The catch-up extension was added after the returns-by-density evidence was
seen, so that evidence cannot test it. The predictions below were not used to
build the model. They are stated before the relevant data were analysed, and
several before the data exist. To make the claim verifiable they should be
deposited in a public registry before they are taken to the data. The
replication package contains the text.

\paragraph{Predictions testable with existing data.}
\begin{enumerate}[label=\textbf{P\arabic*.},leftmargin=*]
\item \emph{Overcrowding is young and fast.} The fall in returns beyond the peak
should be concentrated in clusters whose flow grew faster than their stock, that
is, with a high flow-to-stock ratio $P/S$ (\cref{prop:toofast}). Among mature
clusters the returns curve should be monotone.
\item \emph{Slow clusters die young.} Among clusters born in the same year, those
whose flow grew slowly in their first years should be abandoned more often, and
the hazard should fall sharply above a threshold growth rate
(\cref{prop:toofast}(i)).
\item \emph{Returns follow the stock, density follows the flow.} Holding flow
density fixed, returns should rise with the cluster's depreciated stock of
literature (\cref{prop:mature}(ii): the premium is $(\rho+\xi)T(P,S)$).
\item \emph{Newcomers start slowly where stocks are large and flows are fast.}
The lag between a newcomer's arrival in a cluster and her first paper there, and
the length of her first reference list, should rise with the cluster's stock,
and convexly with its flow as the flow approaches what newcomers can read
(\cref{lem:career}). Incumbents should show no such gradient.
\item \emph{The cross-topic gradient is the catch-up elasticity plus audience.}
The returns gradient from researchers who switch clusters should exceed the log
premium of incumbents over pioneers in the clusters they join, by the gap
between the catch-up elasticity and the premium (\cref{prop:mature}(iv)) and by
the audience effect, which can be measured with citations from outside the
cluster.
\item \emph{Close overlap is pinned; broad neighbourhoods fill.} The volume
elasticity should rise with the similarity radius, from near zero at the overlap
radius toward the fixed-field benchmark at broad radii (Fact~\ref{fact:volume}
already shows this ordering), and the ordering should survive embedding-based
similarity. The ratio of the two elasticities measures how easily topics open
inside existing fields (\cref{prop:nested}).
\item \emph{Founders stay.} The share of a cluster's founding authors who keep
publishing in it, and their returns there relative to their earlier exploratory
work, measure the founders' share of the frontier's value, $\theta_F$. Areas with
a larger $\theta_F$ should gain more from reading tools (\cref{thm:incidence}).
\item \emph{Stepping on toes across areas.} Areas with a larger volume elasticity
of density should show stronger diminishing returns to researchers in citation
output, with $1-\lambda_J^{\rm eq}=\mathcal E|\eps_w|$ area by area
(\cref{prop:sot}).
\item \emph{Search costs and detection.} Google was blocked in mainland China
from 2014 \citep{zhengwang2020}. Relative to Hong Kong-based and overseas
Chinese authors in the same topics, mainland-based authors' close overlaps
should have become less often aware, and their papers more often unaware
duplicates, after 2014 (\cref{lem:credit} with lower $\sigma$). The share of
covering papers they pivot away from cannot be observed, so the prediction is
on the aware share and on the citation penalty of their overlaps.
\end{enumerate}

\paragraph{Predictions for the arXiv rule changes of 2025--2026.}
\begin{enumerate}[label=\textbf{P\arabic*.},leftmargin=*,resume]
\item \emph{Endorsement.} Relative to the other sections, Mathematics in
2026 and all sections in 2026--2027 relative to trend should show fewer pioneer
papers (new clusters) in proportion to the fall in entrants. Local density at
the overlap radius, the flagged overlap rate and the aware share should not
change detectably (\cref{prop:ai}(ii) with small $\mathcal E$).
\item \emph{Review articles.} In computer-science clusters where reviews had been
common, mature density should fall and the number of clusters rise relative to
other computer-science clusters over 2026--2027. Fast-growing clusters there
should overshoot less (higher $g^*$) and newcomers' first papers should come
later and with longer reference lists (\cref{prop:ai}(iv), P3).
\item \emph{Submission cap.} The cap should leave total volume, papers per
author and the overlap rate within a few percent of trend. Submissions by
previously heavy submitters should be replaced partly by submissions in their
co-authors' names.
\item \emph{Reading tools.} When literature-summarising tools are widely used,
high-uptake fields where AI is not the subject should form fewer new clusters,
thicken existing ones, and show more overshooting by fast-growing clusters.
Overlap measured by meaning should rise in dense clusters, and citations per
author-year should rise by much less than the time the tools save, except where
the people who open new clusters use them (\cref{prop:ai}(iii)--(iv),
\cref{thm:incidence}).
\item \emph{Production shocks and closure.} A large, fast rise in production
speed should raise the flow of the densest clusters most, lengthen newcomers'
time to first paper there, and cut entry into them for a while
(\cref{thm:speed}(iv)).
\item \emph{Better search raises the aware share.} If search tools improve, the
share of close overlaps that cite the earlier paper should rise, with no
offsetting fall in the overlap rate in sparse clusters (\cref{prop:ai}(v)).
\end{enumerate}

Two earlier predictions have already been settled. Equal returns (full
dissipation) is rejected, and the catch-up cost replaces it. A detection
channel that runs through the prominence of a paper's announcement is rejected
by the listing-rank design (\cref{sec:causal}), which locates detection in
deliberate search.
